% GENERATED FILE. Edit canonical lessons, then run npm run generate:books.
% canonical-source-hash: fnv1a64:5c40b5959da13339
% canonical-lessons: ZH-C347-zhuyi, ZH-C347-yisi, ZH-W347-si-think, ZH-C347-sikao, ZH-C347-yiyi, ZH-C347-tihui

\chapter{What Do You Mean?}
\label{ch:zh-what-do-you-mean}
\hlchaptermodality{347}

\begin{chapteropening}
\textbf{By the end of this chapter:} \emph{I can ask what something means, say an idea is a good one and say I know from experience.}

Complete the last lesson of chapter 347: I can ask what something means, say an idea is a good one and say I know from experience.
\end{chapteropening}

\section[zhǔyi]{zhǔyi — bright idea; idea of what to do}

\label{lesson:ZH-C347-zhuyi}



\begin{quote}
Before the new one: say the Mandarin for to wrongly assume, then the Mandarin for to make a rough guess.
\end{quote}

\subsection*{You'll want to know: zhǔyi}
\emph{zhǔyi} — \textquotedblleft{}bright idea; idea of what to do\textquotedblright{}.

\begin{quote}
\emph{wǒ juéde zhège zhǔyi hěn hǎo} — I think that's a good idea.
\end{quote}

Its characters come later; for now it is a word for the ear.

\begin{cousinweb}
\emph{zhǔyi} is built from \emph{zhǔ} and \zh{意} 'intent', both inherited from Middle Chinese, \emph{tsyuX} and \emph{'iH} (Guangyun, in Baxter's transcription), which glosses the first 'to be in charge' and the second 'will'. The second is said light here.
\end{cousinweb}

\begin{grammarlens}[title={What you've built}]
One more word for giving reasons.
\end{grammarlens}

\subsection*{Your turn}
\begin{itemize}
\raggedright
  \item \emph{Say it:} \emph{zhǔyi}
  \item \emph{Say it:} \emph{zhǔyi}, once more
  \item \emph{Say it:} \emph{wǒ juéde zhège zhǔyi hěn hǎo} — I think that's a good idea.
  \item \emph{Recall:} say the Mandarin for to wrongly assume, then the Mandarin for to make a rough guess, then say \emph{zhǔyi} again
\end{itemize}

\begin{tcolorbox}[breakable,colback=teal!4,colframe=teal!35!black,title={Before you move on}]
What does zhǔyi mean? (\textquotedblleft{}Bright idea; idea of what to do\textquotedblright{}.) Say it once more.
\end{tcolorbox}
\section[\texorpdfstring{\zh{意思} (yìsi)}{意思 (yìsi)}]{\zh{意思} (yìsi) — meaning; what something means}

\label{lesson:ZH-C347-yisi}



\begin{quote}
Before the new one: say the Mandarin for to make a rough guess, then the Mandarin for a bright idea.
\end{quote}

\subsection*{You'll want to know: \zh{意思}}
\textbf{\zh{意思}} — \emph{yìsi} — \textquotedblleft{}meaning; what something means\textquotedblright{}.

\begin{quote}
\textbf{\zh{这个字是什么意思}} \emph{zhège zì shì shénme yìsi} — What does this character mean?
\end{quote}

The new character, in pieces: \zh{思} — nine strokes in two pieces: on top a box with a cross inside it, five strokes; under it \zh{心} \emph{xīn}, the heart, four strokes. The next lesson puts it in your hand.

\begin{cousinweb}
\textbf{\zh{意思}} (\emph{yìsi}) is built from \zh{意} 'intent' and \zh{思} 'to think', both inherited from Middle Chinese, \emph{'iH} and \emph{si} (Guangyun, in Baxter's transcription), which glosses the first 'will' and the second 'to think of'. The second is said light here.
\end{cousinweb}

\begin{grammarlens}[title={What you've built}]
One more word for giving reasons.
\end{grammarlens}

\subsection*{Your turn}
\begin{itemize}
\raggedright
  \item \emph{Say it:} \emph{yìsi}
  \item \emph{Say it:} \emph{yìsi}, once more
  \item \emph{Say it:} \emph{zhège zì shì shénme yìsi} — What does this character mean?
  \item \emph{Recall:} say the Mandarin for to make a rough guess, then the Mandarin for a bright idea, then say \emph{yìsi} again
\end{itemize}

\begin{tcolorbox}[breakable,colback=teal!4,colframe=teal!35!black,title={Before you move on}]
What does \zh{意思} mean? (\textquotedblleft{}Meaning; what something means\textquotedblright{}.) Say it once more.
\end{tcolorbox}
\section[\texorpdfstring{\zh{思} (sī)}{思 (sī)}]{\zh{思} — a box with a cross, over the heart}

\label{lesson:ZH-W347-si-think}



\begin{quote}
The word first, then the character inside it.

\begin{itemize}
\raggedright
  \item \emph{Recall:} say \emph{yìsi} — \textbf{R1}, one lesson back, before it decays
\end{itemize}

\begin{quote}
\zh{心} \zh{想} \zh{意}
\end{quote}

\emph{Read it:} each one aloud, with nothing said first — \emph{xīn}, \emph{xiǎng}, \emph{yì}
\end{quote}

\subsection*{Script — a box with a cross, over the heart}
\hlblockfigure{\detokenize{figures/ZH-W347-si-think-filmstrip.pdf}}{How \zh{思} is written, stroke by stroke}

\begin{quote}
\zh{思}
\end{quote}

\textbf{Two pieces, top before bottom.} Count the pieces, not the strokes.

\textbf{The top piece, a box with a cross inside, five strokes.} A short vertical down the left side. A horizontal along the top that turns down the right side. A short horizontal across the middle. A vertical down the middle. Then the bottom horizontal that closes the box.

\textbf{The bottom piece, \zh{心}, four strokes}, wider than the box: the left dot, the long curve along the foot that hooks up, the middle dot, the right dot.

5 + 4 = nine, and the last four your hand already knows.

\subsection*{Your turn}
\begin{itemize}
\raggedright
  \item \emph{Write it:} \zh{思}, all nine, with the model in front of you
  \item \emph{Cover:} write it again from memory, and count the pieces, not the strokes
\end{itemize}

\begin{tcolorbox}[breakable,colback=teal!4,colframe=teal!35!black,title={Before you move on}]
How many pieces in \zh{思}? (\textbf{Two}.) Which comes first, the box or the heart? (\textbf{The box, on top}.) How many strokes in the box? (\textbf{Five}.)

Source: \href{https://raw.githubusercontent.com/chanind/hanzi-writer-data/68d10a4b21150cae5e1ebbd223eed289cf32d90c/data/\%E6\%80\%9D.json}{Hanzi Writer Data \zh{思}, pinned stroke paths and medians}.
\end{tcolorbox}
\section[sīkǎo]{sīkǎo — to think hard; to reflect}

\label{lesson:ZH-C347-sikao}



\begin{quote}
Before the new one: say the Mandarin for a bright idea, then the Mandarin for what something means.
\end{quote}

\subsection*{You'll want to know: sīkǎo}
\emph{sīkǎo} — \textquotedblleft{}to think hard; to reflect\textquotedblright{}.

\begin{quote}
\emph{tā xǐhuan yí ge rén sīkǎo} — He likes to think things through alone.
\end{quote}

Its characters come later; for now it is a word for the ear.

\begin{cousinweb}
\emph{sīkǎo} is built from \zh{思} 'to think' and \emph{kǎo} 'to check', both inherited from Middle Chinese, \emph{si} and \emph{khawX} (Guangyun, in Baxter's transcription), which glosses the first 'to think of' and the second 'to check'.
\end{cousinweb}

\begin{grammarlens}[title={What you've built}]
One more word for giving reasons.
\end{grammarlens}

\subsection*{Your turn}
\begin{itemize}
\raggedright
  \item \emph{Say it:} \emph{sīkǎo}
  \item \emph{Say it:} \emph{sīkǎo}, once more
  \item \emph{Say it:} \emph{tā xǐhuan yí ge rén sīkǎo} — He likes to think things through alone.
  \item \emph{Recall:} say the Mandarin for a bright idea, then the Mandarin for what something means, then say \emph{sīkǎo} again
\end{itemize}

\begin{tcolorbox}[breakable,colback=teal!4,colframe=teal!35!black,title={Before you move on}]
What does sīkǎo mean? (\textquotedblleft{}To think hard; to reflect\textquotedblright{}.) Say it once more.
\end{tcolorbox}
\section[yìyì]{yìyì — significance; meaning}

\label{lesson:ZH-C347-yiyi}



\begin{quote}
Before the new one: say the Mandarin for what something means, then the Mandarin for to think hard.
\end{quote}

\subsection*{You'll want to know: yìyì}
\emph{yìyì} — \textquotedblleft{}significance; meaning\textquotedblright{}.

\begin{quote}
\emph{wǒ juéde zhège gōngzuò hěn yǒu yìyì} — I think this work really matters.
\end{quote}

Its characters come later; for now it is a word for the ear.

\begin{cousinweb}
\emph{yìyì} is built from \zh{意} 'intent' and \emph{yì} 'right conduct', both inherited from Middle Chinese, \emph{'iH} and \emph{ngjeH} (Guangyun, in Baxter's transcription), which glosses the first 'will' and the second 'benevolence and right conduct'.
\end{cousinweb}

\begin{grammarlens}[title={What you've built}]
One more word for giving reasons.
\end{grammarlens}

\subsection*{Your turn}
\begin{itemize}
\raggedright
  \item \emph{Say it:} \emph{yìyì}
  \item \emph{Say it:} \emph{yìyì}, once more
  \item \emph{Say it:} \emph{wǒ juéde zhège gōngzuò hěn yǒu yìyì} — I think this work really matters.
  \item \emph{Recall:} say the Mandarin for what something means, then the Mandarin for to think hard, then say \emph{yìyì} again
\end{itemize}

\begin{tcolorbox}[breakable,colback=teal!4,colframe=teal!35!black,title={Before you move on}]
What does yìyì mean? (\textquotedblleft{}Significance; meaning\textquotedblright{}.) Say it once more.
\end{tcolorbox}
\section[\texorpdfstring{\zh{体会} (tǐhuì)}{体会 (tǐhuì)}]{\zh{体会} (tǐhuì) — to know from experience; to come to understand}

\label{lesson:ZH-C347-tihui}



\begin{quote}
Before the new one: say the Mandarin for to think hard, then the Mandarin for the significance of something.
\end{quote}

\subsection*{You'll want to know: \zh{体会}}
\textbf{\zh{体会}} — \emph{tǐhuì} — \textquotedblleft{}to know from experience; to come to understand\textquotedblright{}.

\begin{quote}
\emph{wǒ néng tǐhuì nǐ de xīnqíng} — I know how you feel.
\end{quote}

\begin{cousinweb}
\textbf{\zh{体会}} (\emph{tǐhuì}) is built from \zh{体} 'the body' and \zh{会} 'to meet', both inherited from Middle Chinese, \emph{thejX} and \emph{hwajH} (Guangyun, in Baxter's transcription), which glosses the first 'the body' and the second 'to come together'. How the pair came to mean knowing from experience is not settled.
\end{cousinweb}

\begin{grammarlens}[title={What you've built}]
One more word for giving reasons.
\end{grammarlens}

\subsection*{Your turn}
\begin{itemize}
\raggedright
  \item \emph{Say it:} \emph{tǐhuì}
  \item \emph{Say it:} \emph{tǐhuì}, once more
  \item \emph{Say it:} \emph{wǒ néng tǐhuì nǐ de xīnqíng} — I know how you feel.
  \item \emph{Recall:} say the Mandarin for to think hard, then the Mandarin for the significance of something, then say \emph{tǐhuì} again
\end{itemize}

\begin{tcolorbox}[breakable,colback=teal!4,colframe=teal!35!black,title={Before you move on}]
What does \zh{体会} mean? (\textquotedblleft{}To know from experience; to come to understand\textquotedblright{}.) Say it once more.
\end{tcolorbox}
